\itemhead{p.55 练习 · 第2题（教材p.55）}{XJ-B1-C02-PR-P055-Y1178}
\label{XJ-B1-C02-PR-P055-Y1178}
{\footnotesize 湘教版教材p.55 / PDF第62页 / p.55 练习 / 第2题}
\field{原题}{当 $k$ 为何值时，关于 $x$ 的方程 $x^2+2(k-3)x+4k=0$ 分别满足：
\subq{1}无实数根？
\subq{2}有两正实根？}
\field{参考答案}{（1）$k\in(1,9)$；（2）$k\in(0,1]$，允许判别式为0时的两个相等正根；若额外要求相异，则改为 $(0,1)$。}
{\small \field{关联知识}{Q01 二次函数、零点、判别式与根系关系（教材帮p.65 / PDF74）}
\field{方法与易错点}{按已知零点、顶点或系数选择交点式、顶点式或一般式；根平方和用根系关系消元。 易错：遗漏 $a
e0$；把二重根算成两个不同零点；有两正根只查积正。}
\field{答案依据}{答案册相关题组：PDF第14页（册内页13）；按当前题库原题逐项重算，不按旧版题号机械对套。}}
